%% LyX 2.4.4 created this file.  For more info, see https://www.lyx.org/.
%% Do not edit unless you really know what you are doing.
\documentclass[oneside,english]{book}
\usepackage[T1]{fontenc}
\usepackage[latin9]{inputenc}
\setcounter{secnumdepth}{3}
\setcounter{tocdepth}{3}
\usepackage{units}
\usepackage{amsthm}
\usepackage{amssymb}
\usepackage{geometry}
\geometry{verbose,tmargin=1cm,bmargin=1cm,lmargin=1.5cm,rmargin=1.5cm}
\PassOptionsToPackage{normalem}{ulem}
\usepackage{ulem}

\makeatletter
%%%%%%%%%%%%%%%%%%%%%%%%%%%%%% Textclass specific LaTeX commands.
\theoremstyle{definition}
\newtheorem{example}{\protect\examplename}
\theoremstyle{definition}
\newtheorem{defn}{\protect\definitionname}
\theoremstyle{definition}
\newtheorem{xca}{\protect\exercisename}
\theoremstyle{remark}
\newtheorem{rem}{\protect\remarkname}
\theoremstyle{plain}
\newtheorem{prop}{\protect\propositionname}
\ifx\proof\undefined\
  \newenvironment{proof}[1][\proofname]{\par
    \normalfont\topsep6\p@\@plus6\p@\relax
    \trivlist
    \itemindent\parindent
    \item[\hskip\labelsep
          \scshape
      #1]\ignorespaces
  }{%
    \endtrivlist\@endpefalse
  }
  \providecommand{\proofname}{Proof}
\fi
\theoremstyle{plain}
\newtheorem{lem}{\protect\lemmaname}
\theoremstyle{plain}
\newtheorem{thm}{\protect\theoremname}

%%%%%%%%%%%%%%%%%%%%%%%%%%%%%% User specified LaTeX commands.
\usepackage{hyperref}
\usepackage[usenames,dvipsnames]{xcolor} %adds additional color options
\hypersetup{colorlinks=true, urlcolor=Blue}%Blue

\usepackage{multicol}

\usepackage{tikz}
\usepackage{tikz-3dplot}
\usetikzlibrary{calc,arrows,shapes,backgrounds,matrix,shapes.geometric,fit,trees,positioning}

\usetikzlibrary{patterns}
\usetikzlibrary{plotmarks}
\usepackage{pgfplots}
\pgfplotsset{compat=newest}

\usepackage{calc}
\usepackage[nomessages]{fp}

\usepackage{datetime}

%%%%%commands
\newcommand{\mb}[1]{$\mathbb{#1}$}
\newcommand{\funct}[2]{$f\colon#1\rightarrow#2$}
% Added by lyx2lyx
\setlength{\parskip}{\smallskipamount}
\setlength{\parindent}{0pt}

\makeatother

\usepackage{babel}
\providecommand{\definitionname}{Definition}
\providecommand{\examplename}{Example}
\providecommand{\exercisename}{Exercise}
\providecommand{\lemmaname}{Lemma}
\providecommand{\propositionname}{Proposition}
\providecommand{\remarkname}{Remark}
\providecommand{\theoremname}{Theorem}

\begin{document}

\chapter{Introduction to Sets}

Set theory is a foundation on which (almost) all of mathematics can
be built. 
\begin{itemize}
\item A \emph{set }is a collection of distinct objects called \emph{elements}.
\end{itemize}
This is not really a definition (for a rigorous definition see this
\href{http://en.wikipedia.org/wiki/Set_theory}{link}); to avoid some
messiness we will use this ``intuitive'' (sometimes called ``naive'')
understanding of set and element. Which has good enough for Cantor,
so it should be good enough for us.\footnote{George Cantor founded set theory using the ``naive'' defintions
in 1874 with a single paper. However by 1900 a number of contradictions
were found. Most famously Russell's paradox: ``the set of all sets
that are not members of themselves.'' This set must be a member of
itself, and yet is not a member of itself!} 

\section{Describing Sets}

Lots of new notation will be introduced as we move through the material.
The trade-off for learning all the new notation is certainly worth
it. Think of trying to solve even a simple quadratic, $x^{2}+2x-10=0$,
without the advantage of using $x$ as the unknown and Arabic numbers.
\begin{itemize}
\item Sets are usually denoted with capital letters, e.g., $A,B,C,\dots$
etc; and elements in sets are usually denoted with lowercase letters
$a,b,c,\dots$.
\item If $a$ is an element in the set $A$, we write $a\in A$. If $a$
is not in $A$ we write $a\notin A$ (slashes usually mean negation). 
\end{itemize}
Sets are described in a number of different ways:

Described: 
\begin{example}
$a,b,c$ are elements of the set $A$; also written $a,b,c\in A$.
\end{example}

\begin{example}
$\mathbb{R}$ is the set of all real numbers.
\end{example}

\emph{Tabular }form: 
\begin{example}
$A=\left\{ a,b,c\right\} $, means $a,b,c\in A$.

\begin{itemize}
\item []Dots are used when unlisted elements can be implied:
\end{itemize}
\end{example}

\begin{example}
$A=\left\{ a,b,c,\dots z\right\} $, it is implied that all the letters
in alphabet are in $A$.
\end{example}

\begin{example}
$\mathbb{Z}=\left\{ \dots,-1,0,1,2,\dots\right\} $
\end{example}
\emph{Set builder }form: 
\begin{example}
$\mathbb{Z}^{+}=\left\{ n\in\mathbb{Z}\colon\ensuremath{n>0}\right\} $,
read ``$n$ is an integer \emph{such that }$n$ is greater than $0$.''
\end{example}

\begin{example}
$2\mathbb{Z}=\left\{ n\in\mathbb{Z}\colon n=2m\mbox{ for some }m\in\mathbb{Z}\right\} $.
This set is all the even numbers. So $4\in2\mathbb{Z}$ and $11\notin2\mathbb{Z}$.
\end{example}

\section{Some Basic Definitions}

\emph{A farmer had a 100 feet of fence and wanted to enclose as large
an area as possible. After a little thought, an engineer made a square
resulting in an area of $25^{2}=625\mbox{ ft}^{2}$. A physicist (knowing
he could do better) made a circle resulting in $\pi\left(\nicefrac{100}{2\pi}\right)^{2}\approx795.7\,\mbox{ft}^{2}$.
Finally a mathematician steps forward, breaks off a small piece of
fence, and says ``I define myself to be on the outside!'' (resulting
in an area approximately equal to the area of the Earth). }

How we define things is very important. When trying to prove something,
read the definitions very carefully. For the following definitions,
assume that $A$ and $B$ are both sets. 
\begin{defn}
$A$ is a \emph{subset of $B$ }if every element of $A$ is also an
element of $B$, written $A\subseteq B$. 
\end{defn}
\begin{example}
If $A=\{1,2,3\}$ and $B=\{0,1,2,3\}$then $A\subseteq B$ and $B\not\subseteq A$.
Also note that $A\subseteq A$. 
\end{example}
\begin{itemize}
\item $A\subseteq A$ for any set $A$. Every set is subset of itself. Since
every element of a set $A$ will also be an element of $A$.
\end{itemize}
\begin{example}
$\mathbb{N}\subseteq\mathbb{Z}\subseteq\mathbb{Q}\mathbb{\subseteq R\subseteq C}$:
The natural, integers, rational, real, and complex numbers respectively.
\end{example}
\begin{defn}
The set with no elements is called the \emph{empty set}, denoted $\emptyset$.
\end{defn}
\begin{itemize}
\item $\emptyset=\{x\in A\colon x\notin A\mbox{ for any set }A\}$.
\item Let $B$ be set of numbers that are both irrational and rational real
numbers. $B=\emptyset$.
\item $\emptyset\subseteq A$ for any set $A$, i.e., \emph{the empty set
is a subset of every set}. 
\end{itemize}
\begin{defn}
$A=B$ if and only if $A$ and $B$ have the same elements.
\end{defn}
\begin{itemize}
\item It is often helpful to use the equivalent definition: $A=B$ is $A\subseteq B$
\uline{and} $B\subseteq A$.
\end{itemize}
\begin{example}
Let $A=\{a\colon a=2^{n}\mbox{ for any integer }n\}$ and $B=\{b\colon b=2^{n-2}$
for any integer $n$\}. Show that $A=B$. (Show that $A\subseteq B$)
Let $a\in A$. If $n$ is an integer then there is another integer
$m$ such that $m=n-2$. So then $a=2^{n}=2^{m-2}\in B$. \\
(Show that $B\subseteq A$) Let $b\in B$. If $n$ is an integer then
there is another integer $m$ such that $n-2=m$. So then $b=2^{n-2}=2^{m}\in A$.

\begin{itemize}
\item Note that proving $B\subseteq A$ was nearly identical to proving
$A\subseteq B$ (I literally copied,pasted, and switch a few letters).
When one part of prove is so clearly similar to an already completed
previous part, we will often just write ``similarly it can be shown
that $B\subseteq A$''.
\end{itemize}
\end{example}
\begin{defn}
$A$ is a \emph{proper subset of $B$ }if every element of $A$ is
also an element of $B$ \uline{and} $A\neq B$, written $A\subseteq B$.
\end{defn}
\begin{example}
Let $A=\{a\in\mathbb{Z}\colon a=2^{n}\mbox{ for a positive integer \ensuremath{n}}\}$
and $B=\{a\in\mathbb{Z}\colon a=2n\}$. Show that $A$ is a proper
subset of $B$. 

\begin{itemize}
\item []We must show that $A\subseteq B$ but $A\neq B$. Let $a\in A$.
by definition $a=2^{n}=2\cdot2^{n-1}$ for $n\geq0$. Since $2^{n-1}$
is an integer, $a\in A$. (To show that $A\neq B$ we need only find
1 element in $B$ that is not in $A$. There are many.) $6=2\cdot3\in B$
but $6\notin A$ since $6\neq2^{n}$ for $n$ a positive integer. 
\end{itemize}
\end{example}
\begin{itemize}
\item Our use of$\subseteq$ and $\subset$ is analogous to $\leq$ and
$<$.\footnote{Note that the meaning of set notation varies among authors. $\varsubsetneq$
is often used to denote proper subset, and others use$\subset$ for
subset. }
\end{itemize}
\begin{figure}
\hfill%
\begin{minipage}[t]{0.33\textwidth}%
\begin{center}

\begin{tikzpicture}     	\begin{scope}[scale=.5,shift={(3cm,-5cm)}, fill opacity=0.5,mytext/.style={text opacity=1,font=\large\bfseries}]         
		\draw[fill=red, draw = black] (0,0) circle (2);         
		\draw[fill=green, draw = black] (-0.75,0) circle (1);     
		\node[mytext] at (0,1.5) (B) {\large\textbf{B}};     		
		\node[mytext] at (-1,0) (A) {\large\textbf{A}};     	
	\end{scope}
\end{tikzpicture}
\caption{}{$A\subseteq B$ and $A\subset B$}
\par\end{center}%
\end{minipage}\hfill%
\begin{minipage}[t]{0.33\textwidth}%
\begin{center}

\begin{tikzpicture}     	\begin{scope}[scale=.5,shift={(3cm,-5cm)}, fill opacity=0.5,mytext/.style={text opacity=1,font=\large\bfseries}]         
		\draw[fill=red, draw = black] (0,0) circle (2);         
		\draw[fill=green, draw = black] (0.0,0) circle (2);     
		\node[mytext] at (0,1.5) (B) {\large\textbf{B}};     		
		\node[mytext] at (-1,0) (A) {\large\textbf{A}};     	
	\end{scope}
\end{tikzpicture}
\caption{}{$A\subseteq B$ and $A\not\subset B$}
\par\end{center}%
\end{minipage}\hfill%
\begin{minipage}[t]{0.33\textwidth}%
\begin{center}

\begin{tikzpicture}     	\begin{scope}[scale=.5,shift={(3cm,-5cm)}, fill opacity=0.5,mytext/.style={text opacity=1,font=\large\bfseries}]         
		\draw[fill=red, draw = black] (0,0) circle (2);         
		\draw[fill=green, draw = black] (-2.0,0) circle (1);     
		\node[mytext] at (0,1.5) (A) {\large\textbf{B}};     		
		\node[mytext] at (-2.5,0) (B) {\large\textbf{A}};     	
	\end{scope}
\end{tikzpicture}
\caption{}{$A\not\subset B$}
\par\end{center}%
\end{minipage}\hfill

\caption{Venn diagrams}
\end{figure}

\begin{defn}
Let $A$ be any set. The set of all subsets of $A$ is called the
\emph{power set }of $A$, denoted $\mathcal{P}(A)$. 
\end{defn}
\begin{xca}
Let $A=\{0,1,2,3\}$. Find $\mathcal{P}(A)$ (don't forget $A$ and
$\emptyset$).
\end{xca}
\begin{itemize}
\item Note that: $A\in\mathcal{P}(A)$, $\emptyset\in\mathcal{P}(A)$, and
$\mathcal{P}(\emptyset)=\emptyset$.
\item \label{thm:order-of-power-set}Theorem \ref{thm:-power set order}:
If $A$ be any set with $n$ elements then $\mathcal{P}(A)$ has $2^{n}$
elements. 
\end{itemize}
\begin{rem}
Some summary.

\begin{enumerate}
\item $A=B\iff A\subseteq B$ and $B\subseteq A$.\footnote{``$\iff$'' is equivalent to writing ``if and only if''. Often
this is notated ``iff.''}
\item $A\subseteq A$
\item $\emptyset\subseteq A$
\item $A\in\mathcal{P}(A)$ and $\emptyset\in\mathcal{P}(A)$
\end{enumerate}
\end{rem}

\section{Set Operations}
\begin{defn}
The \emph{union }of sets $A$ and $B$ is defined by $A\cup B=\{x\colon x\in A\mbox{ or }x\in B\}$
\end{defn}
\begin{example}
Let $A=\{1,2,3\}$ and $B=\{2,3,4\}$. $A\cup B=\{1,2,3,4\}$. 
\end{example}
\begin{xca}
$A\cup B=B$ if and only if $A\subseteq B$.
\end{xca}

\begin{example}
$\mathbb{R\cup\mathbb{Z}=\mathbb{R}}$. $\mathbb{Z}\subseteq\mathbb{R}$
so this is true by the above result.
\end{example}

\begin{xca}
$A\cup B=B\cup A$.
\end{xca}
\begin{defn}
The intersection of sets $A$ and $B$ is defined by $A\cap B=\{x\colon x\in A\mbox{ and }x\in B\}$.
\end{defn}
\begin{example}
Let $A=\{1,2,3\}$ and $B=\{2,3,4\}$. $A\cap B=\{2,3\}$.
\end{example}

\begin{example}
Let $A=\{1,2,3\}$ and $B=\{4,5,6\}$. $A\cap B=\emptyset$.
\end{example}

\begin{xca}
$A\cap B=B$ if and only if $B\subseteq A$.
\end{xca}

\begin{xca}
$A\cap B=B\cap A$.
\end{xca}
\begin{figure}[h]
\begin{centering}
\def\firstcircle{(90:1.75cm) circle (2.5cm)} 
\def\secondcircle{(210:1.75cm) circle (2.5cm)} 
\def\thirdcircle{(330:1.75cm) circle (2.5cm)}

\begin{tikzpicture}     
\begin{scope}     
	\fill[red]\firstcircle;       
	\fill[green] \secondcircle;       
	\fill[blue] \thirdcircle;     
\end{scope}     
\begin{scope}
	\clip 
	\secondcircle;       
	\fill[yellow] 
	\firstcircle;     
\end{scope}     
\begin{scope}       
	\clip 
	\secondcircle;       
	\fill[cyan] 
	\thirdcircle;     
\end{scope}     
\begin{scope}       
	\clip 
	\firstcircle;       
	\fill[magenta] 
	\thirdcircle;     
\end{scope}     
\begin{scope}       
	\clip \firstcircle;       
	\clip \secondcircle;       
	\fill[white] \thirdcircle;     
\end{scope}     
\draw 
	\firstcircle node[text=white,above] {$A$};     
	\draw \secondcircle node [text=white,below left] {$B$};     
	\draw \thirdcircle node [text=white,below right] {$C$};     
\begin{scope}[font=\small]       
	\node(0,0) {$A\cap B\cap C$};       
	\draw(30:1.5cm) node {$A\cap C$};       
	\draw(150:1.5cm) node {$A\cap B$};       
\draw(270:1.5cm) node {$B\cap C$};     
\end{scope}   
\end{tikzpicture}\caption{A Venn diagram showing the intersections of three sets, $A$, $B$,
and$C$. }
\par\end{centering}
\end{figure}

\begin{defn}
The \emph{difference of sets }$B$ and $A$ is the set of elements
that are in $B$ and not in $A$; written $B-A$.\footnote{Also called the \emph{relative complement of }$A$ in $B$ usually

denoted $B\setminus A$.}
\end{defn}
\begin{itemize}
\item $B-A=\{x\colon x\in B\mbox{ but }x\notin A\}$ (set builder form).
\end{itemize}
\begin{example}
$\mathbb{Z}-2\mathbb{Z}$ is the set of odd integers.
\end{example}
\begin{defn}
The compliment of a set $A$ is the set of elements (under consideration)
not in $A$; denoted $A^{c}$.\footnote{We are being a little sloppy here by ignoring some complications of

the universal set in trade for an easier introduction to sets. }
\end{defn}
\begin{itemize}
\item $A^{c}=\{x\colon x\notin A\}$.
\end{itemize}
\begin{example}
When considering real numbers, $\mathbb{Z}^{c}=\{x\colon x\notin\mathbb{Z}\}$
is the set of all numbers that are not integers. 
\end{example}
\begin{prop}
$A-B=A\cap B^{c}$.

\begin{proof}
Let $x\in A-B$ so then $x\in A$ and $x\notin B$. By definition
of compliments $x\in B^{c}$. So then $x\in A-B$ and $A-B\subseteq A\cap B^{c}$.

Let $x\in A\cap B^{c}$ so then $x\in A$ and $x\notin B$. By definition
of difference of sets, $x\in A-B$ thus $A\cap B^{c}\subseteq A-B$.

$A-B\subseteq A\cap B^{c}$ and $A\cap B^{c}\subseteq A-B$ imply
$A\cap B^{c}=A-B$.
\end{proof}
\end{prop}
\begin{defn}
The number of elements in a set $A$ is called the \emph{order of
$A$}; denoted $|A|$.\footnote{Also called the \emph{cardinality of $A$.}}
\end{defn}
\begin{itemize}
\item When we write $|A|=n$ it is assumed that $n$ is a non-negative integer.
\end{itemize}
\begin{example}
Let $A=\{a,b,c\}$. $|A|=3$.
\end{example}

\begin{example}
Let $A=\{a,b,c\}$ then $|\mathcal{P}(A)|=2^{3}=8$ by Theorem \ref{thm:order-of-power-set}.
\end{example}

\begin{example}
$|\emptyset|=0$.
\end{example}
\begin{lem}
\label{lem:lemma-for-power-setr-num}If $|A|=n$ and $|B|=n+1$ for
sets $A$ and $B$ then $|\mathcal{P}(B)|=2|\mathcal{P}(A)|$.

\begin{proof}
We can without loss of generality\footnote{Often abbreviated WLOG, ``without loss of generality'' is used before
an assumption in a proof to narrow the premise to some special case
when the proof using the special case can obviously be adapted for
all other cases.} assume that $B=A\cup\{b\}$. So that $A=B-\{a\}$. $B$ has two kinds
of subsets; those containing $b$ and those that do not. Now let $S\in\mathcal{P}(B)$
such that $b\notin S$. Then $S\in\mathcal{P}(A)$, and there are
$n$ such sets. Now let $S'\in\mathcal{P}(B)$ such that $b\in S'$.
Then $S'-\{b\}\subseteq A$, and there are $n$ such sets. As a subsets
of $B$ either contains $b$ or it does not, $|\mathcal{P}(B)|=|\mathcal{P}(A)|+|\mathcal{P}(A)|=2|\mathcal{P}(A)|$.
\end{proof}
\end{lem}
\begin{thm}
\textup{\label{thm:-power set order}$|\mathcal{P}(A)|=2^{n}$ for
any set $A$.}
\end{thm}
\begin{example}
Prove theorem \ref{thm:order-of-power-set}. 
\end{example}
\begin{itemize}
\item []Hint: Use lemma \ref{lem:lemma-for-power-setr-num} and induction.
\end{itemize}

\section{Exercises}
\begin{enumerate}
\item Write in tabular form: 

\begin{enumerate}
\item $A=\{x\colon x\mbox{ is a letter in "do not repeat"}\}$

\begin{enumerate}
\item $3\mathbb{Z}=\left\{ x\in\mathbb{Z}\colon x=3n\mbox{ for some \ensuremath{n\in\mathbb{Z}}}\right\} $
\item $N=\left\{ x\colon x\mbox{ is odd and \ensuremath{x}is even}\right\} $
(recall that $0$ is even)
\end{enumerate}
\item Write these sets in set-builder form:

\begin{enumerate}
\item Let $A=\{3,4,5,6,7,\dots\}$
\item Let $B=\{3,5,7,9,\dots\}$
\item Let $\mathbb{Q}$ be the set of rational numbers. 
\end{enumerate}
\end{enumerate}
\item Find all sets of $A=\{1,2,3\}$.
\item Prove that every nonempty subset has a proper subset. 
\item Prove that if $A$ is a subset of $\emptyset$, then $A=\emptyset$.
\item Prove that $A=\{2,3,4,5\}$ is not a subset of $B=\{x\mid x\mbox{ is even}\}$.
\item Prove that $A$ and $B$ are subsets of $A\cup B$.
\item Prove that $A=A\cup A$ and $A=A\cap A$.
\item Prove that $A\cap B$ is a subset of $A$ and of $B$.
\item Prove that $A\cup B=\emptyset$ implies that $A=\emptyset$ and $B=\emptyset$.
\item Prove that $A\cup B=B$ if and only if $A\subseteq B$.
\item Prove that $A\cap B=B$ if and only if $B\subseteq A$.
\item Prove that $A\cup\left(A\cap B\right)=A$.
\item Prove that $A\cap\left(A\cup B\right)=A$.
\item Prove that $\mathcal{P}(A\cap B)=\mathcal{P}(A)\cap\mathcal{P}(B)$.
\item Prove that $A-B$ is a subset of $A\cup B$.
\item Prove that if $A$ a subset of $B$ then $A\cap B=A$.
\item Let $A\cap B=\emptyset$. Prove that $B\cap A'=B$.
\item Prove that if $A$ is a subset of $B$ then $A\cup B=B$.
\item Prove that $A'-B'=B-A$.
\item Prove that if $A$ is a subset of $B$ then $B'$ is a subset of $A'$.
\item Prove that if $A\cap B=\emptyset$ then $A\cup B'=B'$.
\item Prove that $\left(A\cap B\right)'=A'\cup B'$.
\item Prove that $A\subset B$ implies that $A\cup\left(B-A\right)=B$.
\item $\bigstar$ The \emph{product set }of sets $A$ and $B$ is the set
of all ordered pairs $(a,b)$ where $a\in A$ and $b\in B$; it is
denoted by $A\times B$. Prove that $A\times(B\cap C)=(A\times B)\cap(A\times C)$. 
\end{enumerate}

\end{document}
